\chapter{Differential Equations and the Matrix Exponential}

In the previous chapter, diagonalization uncoupled discrete matrix steps
$\vec{u}_{k+1} = \mat{A}\vec{u}_k$ by finding coordinate axes where the matrix
acts as independent dials. Now we ask the natural continuous question: what
happens when the discrete ticks of a clock become a smooth, continuous flow of
time?

Here is a sentence that sounds like it belongs in a calculus book, not a linear
algebra book: \emph{the rate of change of a quantity is proportional to the
quantity itself.} Out of that one sentence falls the exponential function, the
shape of a cooling cup of coffee, the swing of a frictionless pendulum, the
growth of a population, the decay of a radioactive sample, and the
fundamental structure of linear dynamical systems.

The reason linear algebra cares is that nature rarely tracks just one quantity.
A predator population depends on a prey population, which depends back on the
predator. A mixing tank's salt concentration depends on its neighbor's. A
circuit's voltages and currents push on each other. The position and velocity of
a vibrating mass feed each other. In every one of these, you have a \emph{vector}
of quantities $\vec{u}(t)$, and the rule for how it changes is a \emph{matrix}:
\[
  \frac{d\vec{u}}{dt} = \mat{A}\,\vec{u}.
\]
Eigenvalues and eigenvectors solve this equation directly. Everything you learned
about $\mat{A}\vec{x} = \lambda\vec{x}$ -- which directions $\mat{A}$ leaves
alone, how much it stretches along each -- comes back transformed into the
language of motion through time. The eigenvectors become the natural
\textbf{modes} of the system -- pure, independent motions where every component
scales at the same rate along a fixed line without mixing with the others.
The eigenvalues become growth rates and oscillation frequencies.
Stability -- whether the system settles, blows up, or spins forever -- is read
straight off the signs of the eigenvalues.

Over all of it sits one object, the \textbf{matrix
exponential} $e^{\mat{A}t}$. It is the true matrix generalization of the scalar
$e^{at}$, it packages the whole solution into a single propagator (a linear
transformation that advances any state forward by time $t$), and it is
where diagonalization, calculus, and geometry meet. By the end of the chapter
$e^{\mat{A}t}$ should look less like a formula to memorize and more like the
inevitable answer the scalar case was pointing at all along.

\section{The One-Dimensional Story First}
\label{ch25:scalar}

Before any matrices, nail down the scalar case, because every later
formula is a dressed-up version of it.

Suppose a single number $u(t)$ changes at a rate proportional to its current
size:
\[
  \frac{du}{dt} = a\,u, \qquad u(0) = u_0.
\]
Read it out loud: the bigger $u$ is, the faster it grows (if $a > 0$) or the
faster it shrinks (if $a < 0$). What function has this self-reinforcing property?
You already know it. The exponential. We can even \emph{derive} it instead of
guessing. Separate the variables:
\[
  \frac{du}{u} = a\,dt
  \;\Longrightarrow\;
  \int \frac{du}{u} = \int a\,dt
  \;\Longrightarrow\;
  \ln|u| = a t + C.
\]
Exponentiate both sides and absorb the constant into $u_0$:
\[
  \boxed{\,u(t) = u_0\, e^{a t}.\,}
\]

\begin{intuitionbox}
The number $a$ is a \textbf{growth rate}. If $a > 0$, $u$ runs away to infinity.
If $a < 0$, $u$ decays to zero. If $a = 0$, $u$ just sits at $u_0$ forever. The
sign of $a$ alone determines the long-term behavior, while $|a|$ sets the rate.
In the matrix case, the sign of the real part of each eigenvalue determines
whether its mode grows or decays.
\end{intuitionbox}

There is a second way to write $e^{at}$ that will matter enormously in a moment.
Forget the antiderivative and ask: what power series, when you differentiate it,
reproduces itself times $a$? (For a refresher on Taylor series and term-by-term differentiation, see Appendix~\ref{ch39:calculus}.) Try
\[
  e^{at} = 1 + at + \frac{(at)^2}{2!} + \frac{(at)^3}{3!} + \cdots
         = \sum_{k=0}^{\infty} \frac{(at)^k}{k!}.
\]
Differentiate term by term:
\[
  \frac{d}{dt}\!\left[\sum_{k=0}^\infty \frac{a^k t^k}{k!}\right]
  = \sum_{k=1}^\infty \frac{a^k\, k\, t^{k-1}}{k!}
  = \sum_{k=1}^\infty \frac{a^k t^{k-1}}{(k-1)!}
  = a\sum_{j=0}^\infty \frac{a^j t^{j}}{j!}
  = a\, e^{at}.
\]
So the power series \emph{is} the solution, and it satisfies $\frac{d}{dt}e^{at}
= a\,e^{at}$ with $e^{a\cdot 0} = 1$. Remember this trick. When $a$ becomes a
matrix $\mat{A}$, we cannot ``separate variables'' anymore -- but we \emph{can}
still write down a power series, and it will still differentiate to itself times
$\mat{A}$. That is the matrix exponential, already visible on the horizon.

\begin{focusbox}
Two faces of the same function:
\[
  \frac{du}{dt} = a u,\ u(0)=u_0
  \quad\Longleftrightarrow\quad
  u(t) = u_0 e^{at} = u_0\sum_{k=0}^\infty \frac{(at)^k}{k!}.
\]
The differential equation \emph{defines} the behavior; the exponential
\emph{is} the behavior. Everything in this chapter promotes $a$ to a matrix and
keeps both faces.
\end{focusbox}

\begin{recoverybox}
\textbf{Pause and think (mental check):} If a quantity decays according to
$\frac{du}{dt} = -3u$ with initial size $u(0) = 8$, what is $u(t)$? At what time
does the quantity drop to $1$?
\emph{Check:} $u(t) = 8e^{-3t}$. Setting $8e^{-3t} = 1$ gives $e^{3t} = 8$, so
$t = \frac{\ln 8}{3} = \ln 2 \approx 0.693$.
\end{recoverybox}

\section{Coupled Equations: Why a Matrix Appears}
\label{ch25:coupled}

Now let two quantities interact. The smallest real example makes the point.

\begin{example}[Two tanks trading salt]
\label{ch25:tanks}
Two tanks hold brine. Salt flows between them: tank~1 loses salt to tank~2 and
gains some back. Suppose the amounts $u_1(t)$ and $u_2(t)$ obey
\[
  \frac{du_1}{dt} = -2u_1 + u_2, \qquad
  \frac{du_2}{dt} = \phantom{-}u_1 - 2u_2.
\]
You cannot solve the first equation alone, because $u_2$ is in it; you cannot
solve the second alone, because $u_1$ is in it. They are \textbf{coupled} --
each variable's rate of change depends on the current value of the other, so
neither can be determined in isolation. But stack them into a vector
$\vec{u} = \colvec{u_1\\u_2}$ and the coupling becomes a single matrix:
\[
  \frac{d}{dt}\colvec{u_1\\u_2}
  = \begin{bmatrix} -2 & 1 \\ 1 & -2 \end{bmatrix}\colvec{u_1\\u_2},
  \qquad\text{i.e.}\qquad
  \frac{d\vec{u}}{dt} = \mat{A}\vec{u},
  \quad \mat{A} = \begin{bmatrix} -2 & 1 \\ 1 & -2 \end{bmatrix}.
\]
The derivative of a vector just means ``differentiate each component.'' The
matrix $\mat{A}$ encodes \emph{who pushes on whom}: row $i$ tells you how the
rate of $u_i$ depends on every $u_j$.
\end{example}

This is the general setup. A \textbf{linear system of first-order ordinary
differential equations} with constant coefficients -- ``first-order'' meaning
it involves only first derivatives $\frac{d\vec{u}}{dt}$ (rates of change, not
accelerations), and ``linear'' meaning rates are linear combinations of the
current states -- is
\[
  \frac{d\vec{u}}{dt} = \mat{A}\,\vec{u},
  \qquad \vec{u}(0) = \vec{u}_0,
  \qquad \vec{u}(t)\in\R^n,\ \mat{A}\in\R^{n\times n}.
\]
We want a formula for $\vec{u}(t)$: where is the state at every future (and past)
time, given where it starts?

\begin{intuitionbox}
Think of $\mat{A}\vec{u}$ as a \textbf{velocity field}. At every point $\vec{u}$
of the plane (or of $\R^n$), the matrix hands you an arrow $\mat{A}\vec{u}$ that
says ``if you are here, move \emph{this} way, \emph{this} fast.'' A solution
$\vec{u}(t)$ is a path that, at every instant, flows along the arrow it is
standing on. Solving the system means: drop a particle at $\vec{u}_0$ and let the
arrows carry it. The whole picture is a river, and $\mat{A}$ is its current.
\end{intuitionbox}

\begin{figure}[H]
\centering
\begin{tikzpicture}[scale=1.0,>=Stealth]
  \draw[->,gray!55] (-2.4,0)--(2.4,0) node[right]{$u_1$};
  \draw[->,gray!55] (0,-2.4)--(0,2.4) node[above]{$u_2$};
  % velocity field for A = [[-2,1],[1,-2]]
  \foreach \x in {-2,-1,...,2}{
    \foreach \y in {-2,-1,...,2}{
      \ifnum\x=0
        \ifnum\y=0
        \else
          \pgfmathsetmacro{\vx}{-2*\x + 1*\y}
          \pgfmathsetmacro{\vy}{1*\x - 2*\y}
          \pgfmathsetmacro{\len}{sqrt(\vx*\vx+\vy*\vy)+0.0001}
          \pgfmathsetmacro{\sx}{0.32*\vx/\len}
          \pgfmathsetmacro{\sy}{0.32*\vy/\len}
          \draw[->,sourceblue!70] (\x,\y)--($(\x,\y)+(\sx,\sy)$);
        \fi
      \else
        \pgfmathsetmacro{\vx}{-2*\x + 1*\y}
        \pgfmathsetmacro{\vy}{1*\x - 2*\y}
        \pgfmathsetmacro{\len}{sqrt(\vx*\vx+\vy*\vy)+0.0001}
        \pgfmathsetmacro{\sx}{0.32*\vx/\len}
        \pgfmathsetmacro{\sy}{0.32*\vy/\len}
        \draw[->,sourceblue!70] (\x,\y)--($(\x,\y)+(\sx,\sy)$);
      \fi
    }
  }
  \fill[vecred] (0,0) circle (2.2pt);
  \node[vecred,above right,font=\small] at (0.08,0.08) {$\vec{0}$ (sink)};
\end{tikzpicture}
\caption{The velocity field $\mat{A}\vec{u}$ for the two-tank matrix
$\mat{A}=\protect\begin{bmatrix}-2&1\\1&-2\protect\end{bmatrix}$. Every arrow
points roughly inward: every starting state drifts toward the origin. This is a
\emph{stable} system, and the chapter will explain exactly why from the
eigenvalues.}
\label{ch25:fig-field}
\end{figure}

\section{One Eigenvector at a Time: Decoupling by Modes}
\label{ch25:onemode}

Here is the idea that cracks the whole problem open. We do not know how to solve
$\frac{d\vec{u}}{dt} = \mat{A}\vec{u}$ for a general starting vector. But suppose
we are lucky and the system starts \emph{exactly along an eigenvector}. Then the
matrix has nothing to do but stretch, and the vector equation collapses back into
the scalar equation we already solved.

Let $\mat{A}\vec{x} = \lambda\vec{x}$, with $\vec{x}\neq\vec{0}$. Guess a solution
that stays pointed along $\vec{x}$ for all time, only changing its length:
\[
  \vec{u}(t) = e^{\lambda t}\,\vec{x}.
\]
Test it. The left side:
\[
  \frac{d\vec{u}}{dt} = \frac{d}{dt}\!\left(e^{\lambda t}\right)\vec{x}
  = \lambda\, e^{\lambda t}\,\vec{x}.
\]
The right side:
\[
  \mat{A}\vec{u} = \mat{A}\!\left(e^{\lambda t}\vec{x}\right)
  = e^{\lambda t}\,(\mat{A}\vec{x})
  = e^{\lambda t}\,\lambda\vec{x}
  = \lambda\,e^{\lambda t}\,\vec{x}.
\]
They match. So $\vec{u}(t) = e^{\lambda t}\vec{x}$ is a genuine solution. Along an
eigendirection, the matrix system behaves \emph{exactly} like the one-dimensional
$\frac{du}{dt} = \lambda u$: the eigenvalue $\lambda$ is the growth rate, the
eigenvector $\vec{x}$ is the fixed direction.

\begin{focusbox}
Each eigenpair $(\lambda,\vec{x})$ of $\mat{A}$ gives one pure solution -- a
\textbf{mode} -- of $\frac{d\vec{u}}{dt}=\mat{A}\vec{u}$:
\[
  \vec{u}(t) = e^{\lambda t}\,\vec{x}.
\]
It marches along the fixed line $\spn\{\vec{x}\}$, growing like $e^{\lambda t}$.
The eigenvector sets the \emph{direction}; the eigenvalue sets the \emph{rate}.
\end{focusbox}

\begin{intuitionbox}
This is the deep reason eigenvectors deserve their name (\emph{eigen} =
``own/characteristic''). They are the directions in which the dynamical system
does not mix anything -- it just scales. A general state is a combination of these pure
modes, each evolving at its own rate $e^{\lambda t}$. To solve the system, you
\emph{decompose into modes, let each evolve, and add
them back up.}
\end{intuitionbox}

\subsection{Adding modes: the general solution by superposition}
\label{ch25:superposition}

A single mode describes trajectories that start directly on an eigenvector line.
To solve from an arbitrary initial state, we use the fact that the system is \emph{linear}:
if $\vec{u}(t)$ and $\vec{w}(t)$ are both solutions, then any linear combination
$c_1\vec{u}(t) + c_2\vec{w}(t)$ is also a solution, through the principle of \textbf{superposition}.
(This holds because differentiation and matrix multiplication both distribute over sums:
$\frac{d}{dt}(c_1\vec{u}+c_2\vec{w}) = c_1\mat{A}\vec{u} + c_2\mat{A}\vec{w} = \mat{A}(c_1\vec{u}+c_2\vec{w})$.)

So if $\mat{A}$ ($n\times n$) has $n$ independent eigenvectors $\vec{x}_1,\dots,
\vec{x}_n$ with eigenvalues $\lambda_1,\dots,\lambda_n$, the general solution is
the superposition of all $n$ modes:
\[
  \boxed{\,
  \vec{u}(t) = c_1 e^{\lambda_1 t}\vec{x}_1
             + c_2 e^{\lambda_2 t}\vec{x}_2
             + \cdots
             + c_n e^{\lambda_n t}\vec{x}_n.
  \,}
\]
The constants $c_1,\dots,c_n$ are fixed by the initial condition. Set $t=0$
(every exponential becomes $1$):
\[
  \vec{u}(0) = c_1\vec{x}_1 + \cdots + c_n\vec{x}_n = \vec{u}_0.
\]
That is just ``write $\vec{u}_0$ in the eigenvector basis.'' In matrix form, with
$\mat{X} = [\,\vec{x}_1\ \cdots\ \vec{x}_n\,]$ and $\vec{c}=\colvec{c_1\\\vdots\\c_n}$,
\[
  \mat{X}\vec{c} = \vec{u}_0
  \quad\Longrightarrow\quad
  \vec{c} = \mat{X}\inv\vec{u}_0.
\]

\begin{focusbox}
\textbf{Recipe to solve $\dfrac{d\vec{u}}{dt}=\mat{A}\vec{u},\ \vec{u}(0)=\vec{u}_0$
(diagonalizable case):}
\begin{enumerate}
  \item Find eigenvalues $\lambda_i$ and eigenvectors $\vec{x}_i$ of $\mat{A}$.
  \item Expand the start in that basis: solve $\mat{X}\vec{c}=\vec{u}_0$ for the
        coefficients $c_i$.
  \item Let each coefficient ride its own exponential:
        $\;\vec{u}(t)=\sum_i c_i e^{\lambda_i t}\vec{x}_i.$
\end{enumerate}
\end{focusbox}

\subsection{A fully worked example}
\label{ch25:tanks-solved}

Let us actually solve the two-tank system of Example~\ref{ch25:tanks},
\[
  \mat{A} = \begin{bmatrix} -2 & 1 \\ 1 & -2 \end{bmatrix},
  \qquad
  \vec{u}(0) = \colvec{3\\1}.
\]

\emph{Step 1: eigenvalues.} The characteristic polynomial:
\[
  \det(\mat{A}-\lambda\mat{I})
  = \det\!\begin{bmatrix} -2-\lambda & 1 \\ 1 & -2-\lambda \end{bmatrix}
  = (-2-\lambda)^2 - 1
  = \lambda^2 + 4\lambda + 3
  = (\lambda+1)(\lambda+3).
\]
So $\lambda_1 = -1$ and $\lambda_2 = -3$. Both eigenvalues are negative: because
each mode scales by $e^{\lambda_i t}$, negative rates mean that every component
shrinks toward zero over time. The system is stable, exactly as the inward-pointing
arrows in Figure~\ref{ch25:fig-field} suggested.

\emph{Step 2: eigenvectors.} For $\lambda_1=-1$, solve $(\mat{A}+\mat{I})\vec{x}=0$:
\[
  \begin{bmatrix} -1 & 1 \\ 1 & -1 \end{bmatrix}\vec{x} = \vec{0}
  \;\Longrightarrow\; x_1 = x_2
  \;\Longrightarrow\; \vec{x}_1 = \colvec{1\\1}.
\]
For $\lambda_2=-3$, solve $(\mat{A}+3\mat{I})\vec{x}=0$:
\[
  \begin{bmatrix} 1 & 1 \\ 1 & 1 \end{bmatrix}\vec{x} = \vec{0}
  \;\Longrightarrow\; x_1 = -x_2
  \;\Longrightarrow\; \vec{x}_2 = \colvec{1\\-1}.
\]
(These eigenvectors are perpendicular because $\mat{A}$ is symmetric ($\mat{A}=\mat{A}\T$) ---
a preview of the Spectral Theorem in the very next chapter.)

\emph{Step 3: match the start.} We need $c_1\colvec{1\\1} + c_2\colvec{1\\-1} =
\colvec{3\\1}$, i.e.
\[
  c_1 + c_2 = 3, \qquad c_1 - c_2 = 1
  \;\Longrightarrow\; c_1 = 2,\ c_2 = 1.
\]

\emph{Answer.} Letting each coefficient ride its exponential,
\[
  \vec{u}(t)
  = 2\,e^{-t}\colvec{1\\1} + 1\,e^{-3t}\colvec{1\\-1}
  = \colvec{2e^{-t} + e^{-3t}\\[2pt] 2e^{-t} - e^{-3t}}.
\]
Sanity checks: at $t=0$ this is $\colvec{2+1\\2-1}=\colvec{3\\1}$, correct. As
$t\to\infty$, both exponentials vanish, so $\vec{u}\to\vec{0}$: all the salt
equalizes and drains to the steady state, exactly as the picture promised. The
$e^{-3t}$ part (the imbalance between the tanks) dies fast; the $e^{-t}$ part (the
total amount, along $\colvec{1\\1}$) lingers longest. The \emph{slowest} mode --
the eigenvalue closest to zero -- dominates the late-time behavior
(Figure~\ref{ch25:fig-modal-decomp}).

\begin{figure}[H]
\centering
\begin{tikzpicture}[scale=1.05,>=Stealth]
  \draw[->,gray!50] (-0.6,0)--(3.6,0) node[right,font=\small]{$u_1$};
  \draw[->,gray!50] (0,-1.6)--(0,2.6) node[above,font=\small]{$u_2$};

  \draw[gray!40,dashed,thick] (-0.5,-0.5)--(2.5,2.5) node[above right,font=\small,sourceblue!80!black]{$\spn\{\vec{x}_1\}$ (slow mode, $\lambda_1=-1$)};
  \draw[gray!40,dashed,thick] (-0.5,0.5)--(1.6,-1.6) node[below right,font=\small,vecred!80!black]{$\spn\{\vec{x}_2\}$ (fast mode, $\lambda_2=-3$)};

  \draw[targetgreen!80!black,very thick,->] plot[domain=0:1.8,samples=50,variable=\t]
    ({2*exp(-\t) + exp(-3*\t)}, {2*exp(-\t) - exp(-3*\t)});
  \node[targetgreen!80!black,font=\small,above right] at (2.2,1.2) {trajectory $\vec{u}(t)$};

  \draw[->,sourceblue,thick] (0,0) -- (2,2) node[midway,above left,font=\small]{$2\vec{x}_1$};
  \draw[->,vecred,thick] (0,0) -- (1,-1) node[midway,below left,font=\small]{$1\vec{x}_2$};
  \draw[gray!60,densely dotted] (2,2) -- (3,1);
  \draw[gray!60,densely dotted] (1,-1) -- (3,1);
  \fill[black] (3,1) circle (2pt);
  \node[right,font=\small] at (3.05,1.05) {$\vec{u}(0)=(3,1)\T$};

  \draw[->,sourceblue!70,thick] (0,0) -- (1.213,1.213);
  \draw[->,vecred!70,thick] (0,0) -- (0.2231,-0.2231);
  \draw[gray!45,densely dotted] (1.213,1.213) -- (1.436,0.990);
  \draw[gray!45,densely dotted] (0.2231,-0.2231) -- (1.436,0.990);
  \fill[targetgreen!80!black] (1.436,0.990) circle (1.8pt);
  \node[below right=1pt,font=\footnotesize,targetgreen!80!black] at (1.44,0.95) {$\vec{u}(0.5)$};

  \fill[vecred] (0,0) circle (2pt);
\end{tikzpicture}
\caption{How superposition creates a curved trajectory. At $t=0$, the initial state
$\vec{u}(0)=\protect\colvec{3\\1}$ decomposes into components along the two
eigenvector lines (dashed): $2\vec{x}_1$ (slow mode, blue) and $1\vec{x}_2$ (fast
mode, red). As time runs forward, the fast mode decays rapidly like $e^{-3t}$,
collapsing the red component toward zero. The slow mode decays gently like
$e^{-t}$, keeping the blue component alive. As a result, the sum $\vec{u}(t)$
quickly bends toward the slow eigendirection $\spn\{\vec{x}_1\}$ and follows it into
the origin.}
\label{ch25:fig-modal-decomp}
\end{figure}

\begin{labbox}
Take $\mat{A}=\begin{bmatrix}1&1\\0&2\end{bmatrix}$ and $\vec{u}(0)=\colvec{2\\1}$.
Its eigenvalues are $1$ and $2$ with eigenvectors $\colvec{1\\0}$ and
$\colvec{1\\1}$. Find $c_1,c_2$ (you should get $c_1=1,\ c_2=1$, so both modes are
genuinely present), write $\vec{u}(t)=e^{t}\colvec{1\\0}+e^{2t}\colvec{1\\1}$, and
confirm both components blow up like $e^{2t}$ for large $t$. Which mode wins the
race here -- and why is it the \emph{largest} eigenvalue this time, not the
smallest? (Answer: when eigenvalues are positive, the biggest one dominates
growth; when negative, the one nearest zero dominates the slow decay. Either way
it is the \emph{rightmost} eigenvalue -- largest real part -- that sets the
late-time behavior.)
\end{labbox}

\begin{strangbox}
Strang teaches this as ``$\vec{u}(t) = \sum c_i e^{\lambda_i t}\vec{x}_i$,'' highlighting
how eigenvalues govern dynamical behavior. \emph{Before} Strang, the core geometric insight
is that a system of linear differential equations is solved by identifying the invariant
directions (eigenvectors) and their growth or decay rates (eigenvalues). \emph{After} Strang,
notice this represents diagonalization applied to time evolution, leading directly to the
matrix exponential $e^{\mat{A}t}$.
\end{strangbox}

\section{Stability: Reading the Future from the Eigenvalues}
\label{ch25:stability}

The superposition formula makes the long-term fate of the system transparent.
Each mode carries a factor $e^{\lambda_i t}$. As $t\to\infty$:
\begin{itemize}
  \item if $\lambda_i < 0$, that mode \textbf{decays} to zero;
  \item if $\lambda_i > 0$, that mode \textbf{blows up};
  \item if $\lambda_i = 0$, that mode \textbf{persists} at constant size.
\end{itemize}
The whole solution is a sum of such terms, so its fate is decided by the
\emph{worst} mode -- the one with the largest real part (since, as we are about
to see, eigenvalues may be complex and it is $\operatorname{Re}\lambda$ that sets
growth or decay).

\begin{formalbox}
Consider $\dfrac{d\vec{u}}{dt}=\mat{A}\vec{u}$ with $\mat{A}$ having eigenvalues
$\lambda_1,\dots,\lambda_n$.
\begin{itemize}
  \item The system is \textbf{stable} (every solution $\to\vec{0}$) iff
        $\operatorname{Re}\lambda_i < 0$ for \emph{all} $i$.
  \item It is \textbf{neutrally stable} (solutions stay bounded) iff
        $\operatorname{Re}\lambda_i \le 0$ for all $i$, with the
        zero-real-part eigenvalues not defective (an eigenvalue is \textbf{defective}
        if it lacks a full set of independent eigenvectors, having fewer eigenvectors
        than its algebraic multiplicity).
  \item It is \textbf{unstable} (some solution $\to\infty$) iff
        $\operatorname{Re}\lambda_i > 0$ for some $i$, or some
        eigenvalue with $\operatorname{Re}\lambda_i = 0$ is defective.
\end{itemize}
\end{formalbox}

Why ``real part''? Because eigenvalues can be complex, and that is where
oscillation hides. Recall Euler's formula. If $\lambda = \sigma + i\omega$, then
\[
  e^{\lambda t} = e^{(\sigma + i\omega)t}
  = e^{\sigma t}\,\bigl(\cos\omega t + i\sin\omega t\bigr).
\]
The \emph{real part} $\sigma$ controls the envelope $e^{\sigma t}$ -- growth or
decay -- and the \emph{imaginary part} $\omega$ controls the spinning
$\cos\omega t + i\sin\omega t$ -- oscillation. So:

\begin{intuitionbox}
Split each eigenvalue $\lambda = \sigma + i\omega$ into two jobs:
\begin{center}
\begin{tabular}{ll}
$\sigma = \operatorname{Re}\lambda$ & sets the \textbf{envelope}: $e^{\sigma t}$ grows ($\sigma>0$) or shrinks ($\sigma<0$).\\
$\omega = \operatorname{Im}\lambda$ & sets the \textbf{frequency}: how fast the solution rotates.
\end{tabular}
\end{center}
A pure imaginary eigenvalue ($\sigma=0$) means \emph{undamped oscillation} -- a
frictionless pendulum circling forever. A complex eigenvalue with $\sigma<0$ means
a \emph{damped spiral} into the origin. With $\sigma>0$, a spiral flying out.
\end{intuitionbox}

\begin{figure}[H]
\centering
\begin{tikzpicture}[scale=0.78,>=Stealth]
  % Three phase portraits: sink, spiral sink, center
  % SINK (both real negative)
  \begin{scope}[shift={(0,0)}]
    \draw[->,gray!50] (-1.7,0)--(1.7,0); \draw[->,gray!50] (0,-1.7)--(0,1.7);
    \foreach \a in {0.2,0.55,0.9}{
      \draw[sourceblue,thick,->] plot[domain=0:1.4,samples=30,variable=\t]
        ({\a*1.5*exp(-1.0*\t)},{\a*1.3*exp(-2.2*\t)});
      \draw[sourceblue,thick,->] plot[domain=0:1.4,samples=30,variable=\t]
        ({-\a*1.5*exp(-1.0*\t)},{-\a*1.3*exp(-2.2*\t)});
    }
    \fill[vecred] (0,0) circle (1.6pt);
    \node at (0,-2.2) {\small sink: $\sigma<0$, real};
  \end{scope}
  % SPIRAL SINK (complex, neg real part)
  \begin{scope}[shift={(5,0)}]
    \draw[->,gray!50] (-1.7,0)--(1.7,0); \draw[->,gray!50] (0,-1.7)--(0,1.7);
    \draw[structurepurple,thick,->] plot[domain=0:6.2,samples=120,variable=\t]
      ({1.4*exp(-0.30*\t)*cos(\t r)},{1.4*exp(-0.30*\t)*sin(\t r)});
    \fill[vecred] (0,0) circle (1.6pt);
    \node at (0,-2.2) {\small spiral sink: $\sigma<0$, $\omega\neq0$};
  \end{scope}
  % CENTER (pure imaginary)
  \begin{scope}[shift={(10,0)}]
    \draw[->,gray!50] (-1.7,0)--(1.7,0); \draw[->,gray!50] (0,-1.7)--(0,1.7);
    \foreach \r in {0.5,1.0,1.4}{
      \draw[targetgreen,thick] (0,0) circle (\r);
    }
    \draw[targetgreen,->] (1.4,0)--(1.4,0.25);
    \fill[vecred] (0,0) circle (1.6pt);
    \node at (0,-2.2) {\small center: $\sigma=0$, $\omega\neq0$};
  \end{scope}
\end{tikzpicture}
\caption{Phase portraits keyed to the eigenvalues. Real negative eigenvalues give
a \emph{sink} (paths slide straight in). Complex eigenvalues with negative real
part give a \emph{spiral sink}. Pure imaginary eigenvalues give a \emph{center}
of closed orbits -- perpetual, undamped oscillation. Flip the sign of $\sigma$ and
every arrow reverses: sinks become sources, spirals fly outward.}
\label{ch25:fig-portraits}
\end{figure}

\begin{example}[Planar rotation: complex eigenvalues]
\label{ch25:complex}
Take
\[
  \mat{A} = \begin{bmatrix} 0 & -1 \\ 1 & 0 \end{bmatrix},
  \qquad \frac{d\vec{u}}{dt}=\mat{A}\vec{u}.
\]
Written out, $u_1' = -u_2$ and $u_2' = u_1$. The eigenvalues satisfy
$\lambda^2 + 1 = 0$, so $\lambda = \pm i$: pure imaginary, $\sigma=0$. The theory
predicts a \emph{center} -- closed circular orbits, no growth, no decay. Indeed,
the solution starting at $\colvec{1\\0}$ is
\[
  \vec{u}(t) = \colvec{\cos t\\ \sin t},
\]
which you can verify directly: $\vec{u}'=\colvec{-\sin t\\\cos t}
=\colvec{-u_2\\u_1}=\mat{A}\vec{u}$. The state runs around the unit circle
forever. Notice that $\mat{A}$ is exactly the velocity field whose flow is
rotation -- the matrix that ``generates'' turning. Indeed, as we
will see shortly, $e^{\mat{A}t}$ is precisely the planar rotation matrix through angle $t$.
\end{example}

\begin{trapbox}
\textbf{Sign confusion.} The condition for \emph{stable} continuous dynamics is
$\operatorname{Re}\lambda < 0$, eigenvalues in the \emph{left} half-plane. This is
\emph{not} the same as the rule for discrete iterations $\vec{u}_{k+1} =
\mat{A}\vec{u}_k$, where stability means $|\lambda| < 1$ (eigenvalues inside the
unit \emph{circle}). Continuous time cares about left/right; discrete time cares
about inside/outside the circle. Mixing them up is the classic exam mistake. The
bridge between them is $z = e^{\lambda t}$: the left half-plane maps exactly to
the inside of the unit disk (Figure~\ref{ch25:fig-eigenplane}).
\end{trapbox}

\begin{figure}[H]
\centering
\begin{tikzpicture}[scale=1.0,>=Stealth]
  % left: complex eigenvalue plane, stable = left half
  \begin{scope}
    \fill[sourceblue!12] (-2.1,-1.8) rectangle (0,1.8);
    \draw[->,gray!60] (-2.2,0)--(2.2,0) node[right]{$\operatorname{Re}\lambda$};
    \draw[->,gray!60] (0,-1.9)--(0,1.9) node[above]{$\operatorname{Im}\lambda$};
    \node[sourceblue!70!black,scale=0.8,align=center] at (-1.15,1.45)
      {stable\\(decay)};
    \node[trapred,scale=0.8,align=center] at (1.15,1.45)
      {unstable\\(growth)};
    % sample eigenvalues
    \fill[sourceblue] (-1.4,0.7) circle (1.6pt);
    \fill[sourceblue] (-1.4,-0.7) circle (1.6pt);
    \fill[targetgreen] (0,0.9) circle (1.6pt);
    \fill[targetgreen] (0,-0.9) circle (1.6pt);
    \fill[trapred] (1.0,0.5) circle (1.6pt);
    \node[targetgreen!60!black,scale=0.72,right] at (0.08,0.9) {center};
    \node at (0,-2.4) {\small continuous: stable iff in left half};
  \end{scope}
  % right: discrete eigenvalue plane, stable = inside unit circle
  \begin{scope}[shift={(6.2,0)}]
    \fill[sourceblue!12] (0,0) circle (1.0);
    \draw[->,gray!60] (-2.2,0)--(2.2,0) node[right]{$\operatorname{Re}z$};
    \draw[->,gray!60] (0,-1.9)--(0,1.9) node[above]{$\operatorname{Im}z$};
    \draw[gray!70,thick] (0,0) circle (1.0);
    \fill[sourceblue] (-0.5,0.3) circle (1.6pt);
    \fill[targetgreen] (0.71,0.71) circle (1.6pt);
    \fill[trapred] (1.35,0.5) circle (1.6pt);
    \node[sourceblue!70!black,scale=0.8] at (0,0.0) {$|z|<1$};
    \node at (0,-2.4) {\small discrete: stable iff inside unit circle};
  \end{scope}
  % the bridge
  \draw[->,quotientorange,thick] (2.4,0.0) .. controls (2.9,1.0) and (3.3,1.0) .. (3.8,0.0)
    node[midway,above,quotientorange,scale=0.8] {$z=e^{\lambda t}$};
\end{tikzpicture}
\caption{Two stability regions, and the map between them. \emph{Left:} for
continuous dynamics $\frac{d\vec{u}}{dt}=\mat{A}\vec{u}$, an eigenvalue is
stabilizing when it sits in the \emph{left half-plane} ($\operatorname{Re}\lambda<0$);
on the imaginary axis it gives a neutral center. \emph{Right:} for discrete
iterations $\vec{u}_{k+1}=\mat{A}\vec{u}_k$, the stable region is the \emph{inside
of the unit circle} ($|\lambda|<1$). The exponential map $z=e^{\lambda t}$
carries the left half-plane exactly onto the open unit disk, which is why the two
rules are really one rule seen through $e^{\lambda t}$.}
\label{ch25:fig-eigenplane}
\end{figure}

\begin{walkbox}
A system has eigenvalues $-5$ and $-0.1$. Both decay, so it is stable -- but if
you watched it, what would you \emph{see}? One piece vanishes almost instantly;
the other lingers fifty times longer. The long-term shape of \emph{every}
trajectory is therefore nearly the eigenvector of $-0.1$. Think on your walk: in a
stable system, does the smallest or the largest $|\lambda|$ shape what you actually
observe after a while? (This is exactly why the ``slow mode'' governs the
relaxation time of springs, circuits, and chemical reactions.)
\end{walkbox}

\section{The Matrix Exponential: One Object for the Whole Flow}
\label{ch25:matexp}

We solved the system mode by mode. That is concrete and computational, but it
hides a cleaner truth. In one dimension we never said ``decompose and add'' -- we
just wrote $u(t) = e^{at}u_0$. We want the same compact statement for vectors:
\[
  \vec{u}(t) = \,?\, \cdot \vec{u}_0,
\]
where the ``?'' is a single matrix that takes \emph{any} start $\vec{u}_0$ and
slides it forward to time $t$. That matrix is the \textbf{matrix exponential}
$e^{\mat{A}t}$, and we define it by the very power series we rehearsed in
Section~\ref{ch25:scalar}, with the scalar $a$ replaced by the matrix $\mat{A}$.

\begin{definition}[Matrix exponential]
\label{ch25:def-matexp}
For a square matrix $\mat{M}\in\R^{n\times n}$, define
\[
  e^{\mat{M}}
  \;=\;
  \sum_{k=0}^{\infty}\frac{\mat{M}^k}{k!}
  \;=\;
  \mat{I} + \mat{M} + \frac{\mat{M}^2}{2!} + \frac{\mat{M}^3}{3!} + \cdots .
\]
Here $\mat{M}^0 = \mat{I}$. Just as an infinite series of numbers converges
when its partial sums settle to a definite limit (Appendix~\ref{ch39:calculus}),
an infinite series of matrices converges when the number series in each of its
$n^2$ entries converges. Because the entries of $\mat{M}^k$ grow no faster than
powers of the matrix norm, the series converges entrywise and absolutely for
every square matrix (Appendix~\ref{ch39:integration} develops the calculus of
these matrix series and their term-by-term derivatives). Thus $e^{\mat{M}}$ is
always a well-defined $n\times n$ matrix. In particular,
$e^{\mat{A}t} = \sum_{k=0}^{\infty}\frac{\mat{A}^k t^k}{k!}$.
\end{definition}

\begin{focusbox}
$e^{\mat{A}t}$ is a \emph{matrix}, not a number. It is the ``time-$t$ flow
map'': feed it where you are now, it returns where you will be at time $t$. The
solution of $\frac{d\vec{u}}{dt}=\mat{A}\vec{u}$, $\vec{u}(0)=\vec{u}_0$, is
\[
  \boxed{\ \vec{u}(t) = e^{\mat{A}t}\,\vec{u}_0.\ }
\]
\end{focusbox}

\begin{recoverybox}
\textbf{Pause and think (mental check):} If $\mat{A} = \begin{bmatrix} 2 & 0 \\ 0 & -5 \end{bmatrix}$ is already diagonal, what is $e^{\mat{A}t}$?
\emph{Check:} Because powers of a diagonal matrix are diagonal ($\mat{A}^k = \diag(2^k, (-5)^k)$), each entry computes its own scalar series:
$e^{\mat{A}t} = \diag(e^{2t}, e^{-5t}) = \begin{bmatrix} e^{2t} & 0 \\ 0 & e^{-5t} \end{bmatrix}$.
\end{recoverybox}

Let us prove that boxed claim, and prove it the same way we did for scalars:
differentiate the series term by term and watch it reproduce $\mat{A}$ times
itself.

\begin{theorem}[The matrix exponential solves the system]
\label{ch25:thm-solves}
The matrix-valued function $\Phi(t) = e^{\mat{A}t}$ satisfies
\[
  \frac{d}{dt}\,e^{\mat{A}t} = \mat{A}\,e^{\mat{A}t} = e^{\mat{A}t}\,\mat{A},
  \qquad e^{\mat{A}\cdot 0} = \mat{I}.
\]
Consequently $\vec{u}(t) = e^{\mat{A}t}\vec{u}_0$ solves
$\frac{d\vec{u}}{dt}=\mat{A}\vec{u}$ with $\vec{u}(0)=\vec{u}_0$.
\end{theorem}

\begin{proof}
At $t=0$ every term but the first vanishes, leaving $e^{\mat{A}\cdot 0}=\mat{I}$.
Differentiate the series term by term (legal because it converges absolutely and
uniformly on bounded $t$-intervals; see Appendix~\ref{ch39:integration}):
\[
  \frac{d}{dt}\sum_{k=0}^\infty \frac{\mat{A}^k t^k}{k!}
  = \sum_{k=1}^\infty \frac{\mat{A}^k\, k\, t^{k-1}}{k!}
  = \sum_{k=1}^\infty \frac{\mat{A}^k t^{k-1}}{(k-1)!}.
\]
Factor one $\mat{A}$ out of the front (it commutes with every $\mat{A}^{k-1}$) and
reindex $j=k-1$:
\[
  = \mat{A}\sum_{j=0}^\infty \frac{\mat{A}^j t^j}{j!}
  = \mat{A}\,e^{\mat{A}t}.
\]
Pulling the $\mat{A}$ out on the right instead gives $e^{\mat{A}t}\mat{A}$; the
two agree because $\mat{A}$ commutes with its own powers. Finally, with
$\vec{u}(t) = e^{\mat{A}t}\vec{u}_0$,
\[
  \frac{d\vec{u}}{dt} = \Bigl(\tfrac{d}{dt}e^{\mat{A}t}\Bigr)\vec{u}_0
  = \mat{A}e^{\mat{A}t}\vec{u}_0 = \mat{A}\vec{u}(t),
\]
and $\vec{u}(0)=\mat{I}\vec{u}_0=\vec{u}_0$.
\end{proof}

\begin{connectionbox}
This is the same proof, letter for letter, as the calculus proof that
$\frac{d}{dt}e^{at} = a\,e^{at}$. We replaced the scalar $a$ by the
matrix $\mat{A}$ and changed nothing else. A scalar identity, when evaluated on
matrices, becomes an operator identity describing entire dynamic systems.
Calculus supplies the power series and term-by-term differentiation
(Appendix~\ref{ch39:integration}); linear algebra supplies the algebra of
matrix powers $\mat{A}^k$. The matrix exponential connects the two subjects.
\end{connectionbox}

\subsection{Properties, and the one trap}
\label{ch25:props}

The matrix exponential keeps most of the algebra you expect, with one sharp
exception you must respect.

\begin{proposition}[Algebra of $e^{\mat{A}}$]
\label{ch25:prop-algebra}
For square $\mat{A},\mat{B}$ of the same size:
\begin{enumerate}
  \item $e^{\mat{0}} = \mat{I}$.
  \item $\bigl(e^{\mat{A}t}\bigr)\inv = e^{-\mat{A}t}$; in particular
        $e^{\mat{A}t}$ is \emph{always invertible}.
  \item $e^{\mat{A}(s+t)} = e^{\mat{A}s}\,e^{\mat{A}t}$ (the flow composes:
        evolving $s$ then $t$ equals evolving $s+t$).
  \item \textbf{If $\mat{A}\mat{B}=\mat{B}\mat{A}$,} then
        $e^{\mat{A}+\mat{B}} = e^{\mat{A}}e^{\mat{B}}$.
        If they do \emph{not} commute, this is generally \emph{false}.
\end{enumerate}
\end{proposition}

Property~(2) is worth a pause. If $\mat{A}$ has eigenvalues $\lambda_1,\dots,\lambda_n$,
the eigenvalues of $e^{\mat{A}t}$ are $e^{\lambda_1 t},\dots,e^{\lambda_n t}$.
Because the determinant of any matrix equals the product of its eigenvalues,
\[
  \det(e^{\mat{A}t}) = \prod_{i=1}^n e^{\lambda_i t} = e^{t\sum_{i=1}^n \lambda_i} = e^{t\,\tr\mat{A}},
\]
where $\tr\mat{A}$, the \textbf{trace}, is the sum of the entries on the main
diagonal (which always equals the sum of all eigenvalues). Because the scalar
exponential $e^{t\,\tr\mat{A}}$ is never zero, $\det(e^{\mat{A}t}) \neq 0$ for all
$t$. The flow map is always invertible: it can never flatten space into a
lower-dimensional slice, two distinct initial states can never collide, and
trajectories never cross. That is the matrix face of the uniqueness of solutions
to differential equations.

\begin{trapbox}
\textbf{$e^{\mat{A}+\mat{B}} \neq e^{\mat{A}}e^{\mat{B}}$ when $\mat{A}$ and
$\mat{B}$ do not commute.} The familiar law $e^{x+y}=e^xe^y$ relies on
$xy=yx$, which is automatic for numbers but not for matrices. The cleanest way to
see the failure: expand both sides to second order. The cross term in
$e^{\mat{A}+\mat{B}}$ is $\tfrac12(\mat{A}+\mat{B})^2 = \tfrac12(\mat{A}^2 +
\mat{A}\mat{B} + \mat{B}\mat{A} + \mat{B}^2)$, while in $e^{\mat{A}}e^{\mat{B}}$ it
is $\tfrac12\mat{A}^2 + \mat{A}\mat{B} + \tfrac12\mat{B}^2$. They differ by
$\tfrac12(\mat{A}\mat{B}-\mat{B}\mat{A})$ -- the \textbf{commutator}, which
measures the failure of two operators to commute. Without commutativity, the
exponential of a sum does not factor into a product of exponentials -- an
obstruction studied systematically in Lie theory.
\end{trapbox}

\section{Computing $e^{\mat{A}t}$ by Diagonalization}
\label{ch25:diag}

The power-series definition is foundational, but difficult to evaluate directly.
Diagonalization provides a practical computational route, and reveals that the matrix
exponential unifies the mode expansion from Section~\ref{ch25:superposition} into a single
operator.

Suppose $\mat{A}$ is diagonalizable: $\mat{A} = \mat{X}\mat{\Lambda}\mat{X}\inv$,
where $\mat{\Lambda} = \diag(\lambda_1,\dots,\lambda_n)$ holds the eigenvalues and
$\mat{X}$ holds the eigenvectors as columns. The key observation is that powers
of $\mat{A}$ telescope:
\[
  \mat{A}^k
  = (\mat{X}\mat{\Lambda}\mat{X}\inv)^k
  = \mat{X}\mat{\Lambda}^k\mat{X}\inv,
\]
because the inner $\mat{X}\inv\mat{X}$ pairs cancel. Feed this into the series:
\[
  e^{\mat{A}t}
  = \sum_{k=0}^\infty \frac{\mat{A}^k t^k}{k!}
  = \sum_{k=0}^\infty \frac{\mat{X}\mat{\Lambda}^k\mat{X}\inv t^k}{k!}
  = \mat{X}\!\left(\sum_{k=0}^\infty \frac{\mat{\Lambda}^k t^k}{k!}\right)\!\mat{X}\inv
  = \mat{X}\, e^{\mat{\Lambda}t}\,\mat{X}\inv.
\]
And a diagonal matrix exponentiates entrywise -- because powers of a diagonal
matrix are just the powers of its entries:
\[
  e^{\mat{\Lambda}t}
  = \diag\!\bigl(e^{\lambda_1 t}, e^{\lambda_2 t}, \dots, e^{\lambda_n t}\bigr).
\]

\begin{focusbox}
\textbf{The master formula.} If $\mat{A}=\mat{X}\mat{\Lambda}\mat{X}\inv$, then
\[
  e^{\mat{A}t}
  = \mat{X}
    \begin{bmatrix}
      e^{\lambda_1 t} & & \\
      & \ddots & \\
      & & e^{\lambda_n t}
    \end{bmatrix}
    \mat{X}\inv.
\]
Read it as a sentence: \emph{change to eigen-coordinates ($\mat{X}\inv$),
exponentiate each one-dimensional flow ($e^{\lambda_i t}$), change back
($\mat{X}$).} The hard matrix problem became $n$ trivial scalar problems.
\end{focusbox}

\begin{figure}[H]
\centering
\begin{tikzpicture}[scale=1.0,>=Stealth,
    box/.style={draw,rounded corners=2pt,minimum width=2.2cm,minimum height=0.95cm,align=center,font=\small}]
  \node[box,fill=sourceblue!8] (u0) at (0,1.6) {start $\vec{u}_0$\\(real coords)};
  \node[box,fill=sourceblue!8] (ut) at (6.4,1.6) {answer $\vec{u}(t)$\\(real coords)};
  \node[box,fill=targetgreen!10] (c) at (0,-1.0) {$\vec{c}=\mat{X}\inv\vec{u}_0$\\(eigen-coords)};
  \node[box,fill=targetgreen!10] (ct) at (6.4,-1.0) {$e^{\mat{\Lambda}t}\vec{c}$\\(eigen-coords)};
  \draw[->,gray!70,dashed] (u0) -- node[above,font=\small]{$e^{\mat{A}t}$ (hard, coupled)} (ut);
  \draw[->,quotientorange,thick] (u0) -- node[left,font=\small]{$\mat{X}\inv$} (c);
  \draw[->,structurepurple,thick] (c) -- node[below,align=center,font=\small]{multiply each entry by $e^{\lambda_i t}$\\(easy, uncoupled)} (ct);
  \draw[->,quotientorange,thick] (ct) -- node[right,font=\small]{$\mat{X}$} (ut);
\end{tikzpicture}
\caption{The master formula as a detour. Going straight across the top
($e^{\mat{A}t}$) means solving a coupled system. Instead drop down into
eigen-coordinates with $\mat{X}\inv$, where the system splits into $n$
independent scalar flows, scale each by its own $e^{\lambda_i t}$, then ride
$\mat{X}$ back up. The two routes give the same $\vec{u}(t)$ -- that equality
\emph{is} $e^{\mat{A}t}=\mat{X}\,e^{\mat{\Lambda}t}\,\mat{X}\inv$.}
\label{ch25:fig-detour}
\end{figure}

\begin{intuitionbox}
This represents diagonalization applied to continuous time evolution (Figure~\ref{ch25:fig-detour}).
In the eigenbasis, the coupled system uncouples completely: each eigen-coordinate
$w_i = (\mat{X}\inv\vec{u})_i$ obeys its own scalar differential equation $w_i' = \lambda_i w_i$, so
$w_i(t) = e^{\lambda_i t} w_i(0)$. The matrices $\mat{X}\inv$ and $\mat{X}$ convert
vectors between the standard coordinate frame and the decoupled eigenbasis.
\end{intuitionbox}

And the connection back to modes is immediate. Multiply the master formula by
$\vec{u}_0$ and write $\mat{X}\inv\vec{u}_0 = \vec{c}$:
\[
  \vec{u}(t) = e^{\mat{A}t}\vec{u}_0
  = \mat{X}\,e^{\mat{\Lambda}t}\,\vec{c}
  = \sum_{i=1}^n c_i\, e^{\lambda_i t}\,\vec{x}_i.
\]
Multiplying $e^{\mat{A}t}$ by $\vec{u}_0$ recovers the superposition formula of
Section~\ref{ch25:superposition} directly.

\begin{example}[Exponentiating the two-tank matrix]
\label{ch25:tanks-exp}
For $\mat{A} = \begin{bmatrix}-2&1\\1&-2\end{bmatrix}$ we found
$\lambda_1=-1,\lambda_2=-3$ with $\vec{x}_1=\colvec{1\\1},\vec{x}_2=\colvec{1\\-1}$.
So
\[
  \mat{X} = \begin{bmatrix}1&1\\1&-1\end{bmatrix},
  \qquad
  \mat{X}\inv = \frac12\begin{bmatrix}1&1\\1&-1\end{bmatrix}.
\]
Then
\[
  e^{\mat{A}t}
  = \mat{X}
    \begin{bmatrix}e^{-t}&0\\0&e^{-3t}\end{bmatrix}
    \mat{X}\inv
  = \frac12
    \begin{bmatrix}
      e^{-t}+e^{-3t} & e^{-t}-e^{-3t}\\[2pt]
      e^{-t}-e^{-3t} & e^{-t}+e^{-3t}
    \end{bmatrix}.
\]
Apply it to $\vec{u}_0=\colvec{3\\1}$:
\[
  e^{\mat{A}t}\colvec{3\\1}
  = \frac12\colvec{3(e^{-t}+e^{-3t}) + (e^{-t}-e^{-3t})\\[2pt]
                   3(e^{-t}-e^{-3t}) + (e^{-t}+e^{-3t})}
  = \colvec{2e^{-t}+e^{-3t}\\[2pt] 2e^{-t}-e^{-3t}},
\]
matching Section~\ref{ch25:tanks-solved} exactly. The single matrix $e^{\mat{A}t}$
carries the whole solution for \emph{every} starting vector at once.
\end{example}

\begin{labbox}
Compute $e^{\mat{A}t}$ by hand for the rotation generator
$\mat{A}=\begin{bmatrix}0&-1\\1&0\end{bmatrix}$ of Example~\ref{ch25:complex}, but
\emph{without} diagonalizing -- use the series. Show $\mat{A}^2=-\mat{I}$, hence
$\mat{A}^3=-\mat{A}$, $\mat{A}^4=\mat{I}$, and the powers cycle. Group the series
into even and odd terms and recognize the Taylor series of cosine and sine to get
\[
  e^{\mat{A}t} = \begin{bmatrix}\cos t & -\sin t\\ \sin t & \cos t\end{bmatrix}.
\]
The matrix exponential of the ``turning'' matrix is the rotation by angle $t$.
Beautiful, and a perfect mirror of $e^{i t} = \cos t + i\sin t$.
\end{labbox}

\begin{aestheticbox}
The scalar identity $e^{it}=\cos t + i\sin t$
-- Euler's formula -- is the
counterpart of the matrix statement ``the exponential of the rotation generator is the
rotation matrix.'' The imaginary unit $i$ corresponds algebraically to the matrix
$\begin{bmatrix}0&-1\\1&0\end{bmatrix}$; squaring it gives $-\mat{I}$,
just as $i^2=-1$. Complex numbers, rotations, oscillation, and the matrix
exponential are interconnected facets of the same underlying structure.
\end{aestheticbox}

\section{When You Cannot Diagonalize}
\label{ch25:defective}

The master formula assumed $n$ independent eigenvectors. What if $\mat{A}$ is
\textbf{defective} -- lacking a complete set of $n$ independent eigenvectors, like
$\mat{A}=\begin{bmatrix}\lambda&1\\0&\lambda\end{bmatrix}$? Although the eigenvector
decomposition is incomplete, the matrix exponential remains well-defined: its defining
power series converges for \emph{every} square matrix. You can evaluate it directly by
splitting $\mat{A}$ into a scalar multiple of $\mat{I}$ plus a \textbf{nilpotent} remainder
-- a matrix that vanishes entirely when raised to a sufficiently high power.

Write $\mat{A} = \lambda\mat{I} + \mat{N}$ with $\mat{N}=\begin{bmatrix}0&1\\0&0
\end{bmatrix}$. These commute (every matrix commutes with scalar multiples of the
identity $\lambda\mat{I}$), so by Proposition~\ref{ch25:prop-algebra}(4),
\[
  e^{\mat{A}t} = e^{\lambda t\mat{I}}\, e^{\mat{N}t}
  = e^{\lambda t}\, e^{\mat{N}t}.
\]
Now $\mat{N}^2 = \mat{0}$ (so $\mat{N}$ is nilpotent of degree 2), so its series
truncates cleanly after two terms:
\[
  e^{\mat{N}t} = \mat{I} + \mat{N}t = \begin{bmatrix}1 & t\\ 0 & 1\end{bmatrix}.
\]
Therefore
\[
  e^{\mat{A}t} = e^{\lambda t}\begin{bmatrix}1 & t\\ 0 & 1\end{bmatrix}
  = \begin{bmatrix}e^{\lambda t} & t\,e^{\lambda t}\\ 0 & e^{\lambda t}\end{bmatrix}.
\]

\begin{recoverybox}
\textbf{Pause and think (mental check):} If $\mat{N} = \begin{bmatrix} 0 & 3 \\ 0 & 0 \end{bmatrix}$, what is $\mat{N}^2$, and what is $e^{\mat{N}t}$?
\emph{Check:} $\mat{N}^2 = \begin{bmatrix}0&0\\0&0\end{bmatrix} = \mat{0}$, so only two terms survive the power series:
$e^{\mat{N}t} = \mat{I} + \mat{N}t = \begin{bmatrix} 1 & 3t \\ 0 & 1 \end{bmatrix}$.
\end{recoverybox}

Notice the brand-new ingredient: a factor of $t$ multiplying an exponential. A
defective eigenvalue produces solutions involving $t\,e^{\lambda t}$ alongside
$e^{\lambda t}$. Geometrically, the deficient direction leaks polynomial
growth on top of the exponential.

\begin{remark}
The general structure behind this is the \textbf{Jordan canonical form}
(introduced in Chapter~\ref{ch28:jordan}) -- every matrix can be expressed as
$\mat{X}(\mat{\Lambda}+\mat{N})\mat{X}\inv$ with $\mat{\Lambda}$ diagonal,
$\mat{N}$ nilpotent, and $\mat{\Lambda}\mat{N}=\mat{N}\mat{\Lambda}$. Because
the pieces commute, $e^{\mat{A}t}$ always decomposes into pure exponentials
multiplied by polynomials in $t$. For now, keep the headline: the flow map
$e^{\mat{A}t}$ \emph{exists and is computable for every matrix $\mat{A}$},
diagonalizable or not.
\end{remark}

\begin{connectionbox}
Higher-order scalar ODEs can be rewritten directly as first-order matrix systems, and this is where the
$t\,e^{\lambda t}$ terms you saw in calculus come from. The equation
$y'' + 2y' + y = 0$ has the repeated root $\lambda=-1$, and its solutions are
$e^{-t}$ \emph{and} $t\,e^{-t}$ -- exactly the defective pattern above. Setting
$\vec{u}=\colvec{y\\y'}$ turns the second-order equation into
$\vec{u}'=\begin{bmatrix}0&1\\-1&-2\end{bmatrix}\vec{u}$, a defective $2\times2$
system. The standard rule of multiplying by $t$ for repeated roots in differential
equations is the scalar manifestation of a nilpotent operator.
\end{connectionbox}

\section{A Modern Application}
\label{ch25:application}

\begin{applicationbox}
\textbf{Continuous-time Markov chains and the engines behind search,
epidemics, and chemistry.} Many real systems hop between states at random rates:
which web page a surfer visits next, whether a person is Susceptible /
Infected / Recovered, which conformation a protein is in. The probability
distribution $\vec{p}(t)$ over states evolves by a linear ODE,
\[
  \frac{d\vec{p}}{dt} = \mat{Q}\T\vec{p},
\]
where $\mat{Q}$ is the \emph{generator} (rate) matrix: off-diagonal $Q_{ij}\ge0$
is the rate of jumping $i\to j$, and each row sums to zero so probability is
conserved. The solution is exactly the matrix exponential,
\[
  \vec{p}(t) = e^{\mat{Q}\T t}\,\vec{p}(0).
\]
The eigenvalue $\lambda=0$ (every generator has one, because rows sum to zero)
carries the \emph{steady-state distribution} -- the mode that never decays, the
long-run behavior of the chain. Every other eigenvalue has negative real part, and
the one closest to zero sets the \emph{mixing time}: how long until the chain
forgets where it started. Search ranking, disease forecasting, queueing theory,
and molecular kinetics all run on $e^{\mat{Q}t}$. The signs and magnitudes of
$\mat{Q}$'s eigenvalues determine whether the chain settles to a steady state
and how quickly it mixes.
\end{applicationbox}

\begin{applicationbox}
\textbf{The exponential of a network.} In graph theory and neuroscience, if
$\mat{A}$ is the adjacency matrix of a network, then $(e^{\mat{A}})_{ij}$ -- the
$ij$ entry of the matrix exponential -- is a weighted count of \emph{all} walks
from node $i$ to node $j$, with longer walks down-weighted by $1/k!$. This
``communicability'' is a standard centrality measure: it scores how easily
influence (signals, rumors, current) spreads between nodes. Because
$e^{\mat{A}}=\sum \mat{A}^k/k!$ and $(\mat{A}^k)_{ij}$ counts length-$k$ walks, the
matrix exponential of a graph is, quite literally, a generating function for its
connectivity.
\end{applicationbox}

\section{The Geometry of the Flow}
\label{ch25:geometry}

The geometry is the point, so it is worth collecting the pictures in one place. The
matrix $\mat{A}$ defines a velocity field; the family $\{e^{\mat{A}t}\}_{t\in\R}$
is the \emph{flow} it generates -- a one-parameter family of matrices, starting
at $\mat{I}$ when $t=0$, sliding every point of space along the river as $t$
increases.

\begin{figure}[H]
\centering
\begin{tikzpicture}[scale=1.05,>=Stealth]
  % node1: sink with two straight eigendirections (the two-tank system)
  \begin{scope}
    \draw[->,gray!50] (-2.2,0)--(2.2,0) node[right]{$u_1$};
    \draw[->,gray!50] (0,-2.2)--(0,2.2) node[above]{$u_2$};
    % eigen-lines
    \draw[gray!45,dashed] (-2,-2)--(2,2);
    \draw[gray!45,dashed] (-2,2)--(2,-2);
    \node[gray!60,scale=0.8] at (1.9,1.7) {$\vec{x}_1$};
    \node[gray!60,scale=0.8] at (-1.9,1.7) {$\vec{x}_2$};
    % a few trajectories bending toward slow eigendirection x1
    \foreach \s in {-1,1}{
      \draw[sourceblue,thick,->] plot[domain=0:1.5,samples=40,variable=\t]
        ({\s*(1.0*exp(-1*\t) + 1.4*exp(-3*\t))},
         {\s*(1.0*exp(-1*\t) - 1.4*exp(-3*\t))});
      \draw[sourceblue,thick,->] plot[domain=0:1.5,samples=40,variable=\t]
        ({\s*(1.4*exp(-1*\t) - 1.0*exp(-3*\t))},
         {\s*(1.4*exp(-1*\t) + 1.0*exp(-3*\t))});
    }
    \fill[vecred] (0,0) circle (1.8pt);
    \node at (0,-2.6) {\small stable node: $\lambda=-1,-3$};
  \end{scope}
  % node2: saddle (one positive, one negative eigenvalue)
  \begin{scope}[shift={(6.2,0)}]
    \draw[->,gray!50] (-2.2,0)--(2.2,0) node[right]{$u_1$};
    \draw[->,gray!50] (0,-2.2)--(0,2.2) node[above]{$u_2$};
    % stable manifold (vertical), unstable (horizontal)
    \draw[targetgreen,thick,>-<] (0,-2)--(0,2);
    \draw[vecred,thick,<->] (-2,0)--(2,0);
    \node[targetgreen!80!black,scale=0.75,left] at (-0.08,1.7) {$\lambda = -1$ (stable)};
    \node[vecred!80!black,scale=0.75,below] at (1.2,-0.08) {$\lambda = +1$ (unstable)};
    \draw[structurepurple,thick,->] plot[domain=-1.3:1.3,samples=40,variable=\t]
      ({0.5*exp(\t)},{0.5*exp(-\t)});
    \draw[structurepurple,thick,->] plot[domain=-1.3:1.3,samples=40,variable=\t]
      ({-0.5*exp(\t)},{-0.5*exp(-\t)});
    \draw[structurepurple,thick,->] plot[domain=-0.75:0.75,samples=40,variable=\t]
      ({1.0*exp(\t)},{1.0*exp(-\t)});
    \draw[structurepurple,thick,->] plot[domain=-0.75:0.75,samples=40,variable=\t]
      ({-1.0*exp(\t)},{-1.0*exp(-\t)});
    \fill[vecred] (0,0) circle (1.8pt);
    \node at (0,-2.6) {\small saddle: $\lambda=+1,-1$};
  \end{scope}
\end{tikzpicture}
\caption{Two phase portraits. \emph{Left:} both eigenvalues negative -- a stable
node. Every trajectory slides into the origin, hugging the slow eigendirection
$\vec{x}_1$ (eigenvalue $-1$) on its way in, because the fast mode ($-3$) dies
first. \emph{Right:} one positive, one negative eigenvalue -- a \emph{saddle}.
Along the green axis (stable, $\lambda=-1$) points fall in; along the red axis
(unstable, $\lambda=+1$) they fly out; generic trajectories swoop in, turn, and
escape. A single positive eigenvalue is enough to make a system unstable.}
\label{ch25:fig-node-saddle}
\end{figure}

\begin{intuitionbox}
The full zoo of $2\times2$ behaviors is organized entirely by the two
eigenvalues:
\begin{itemize}
  \item both real, same sign $\to$ \textbf{node} (in if negative, out if positive);
  \item real, opposite signs $\to$ \textbf{saddle} (always unstable);
  \item complex, nonzero real part $\to$ \textbf{spiral} (in if $\sigma<0$);
  \item pure imaginary $\to$ \textbf{center} (closed orbits, neutrally stable).
\end{itemize}
Two numbers -- the eigenvalues -- and you know the qualitative shape of every
trajectory before computing a single point. That is the power the spectrum gives
you over a dynamical system.
\end{intuitionbox}

\begin{strangbox}
Strang frames the climax as: \emph{the eigenvalues of $\mat{A}$ control
$e^{\mat{A}t}$}, and the four fundamental questions (does it grow, decay,
oscillate, blow up?) are answered by the locations of the $\lambda$'s. \emph{After}
Strang, the natural next step is to ask what happens when $\mat{A}$ is symmetric
(real eigenvalues, orthogonal eigenvectors, no spirals -- the spectral theorem,
next chapter) and what $e^{\mat{A}t}$ means when $\mat{A}$ is skew-symmetric
(pure rotation, energy conserved). The matrix exponential is the doorway from
static linear algebra into the dynamics of the world.
\end{strangbox}

\begin{historybox}
For two centuries after Newton, the grand goal of mathematical physics was discovering explicit algebraic formulas for physical paths and planetary orbits. But in the 1880s, the French mathematician Henri Poincar\'e confronted the \textbf{three-body problem} in celestial mechanics -- predicting the gravitational motions of three mutually attracting bodies, such as the Sun, Earth, and Moon -- and discovered that exact closed-form algebraic formulas in terms of elementary functions are impossible. Armed only with pencil-and-paper calculus, coordinate transformations, and power series, Poincar\'e revolutionized mathematics by shifting from algebraic formulas to qualitative geometry. Instead of hunting for an elusive formula for $\vec{u}(t)$, he asked what the global geometry of all trajectories looks like in space. He invented the \textbf{phase portrait} -- a geometric atlas displaying all possible trajectory curves simultaneously -- showing that a linear system $\frac{d\vec{u}}{dt} = \mat{A}\vec{u}$ generates a directional velocity field whose flow lines swirl, focus, saddle, or spiral around equilibrium points. Decades later, Rudolf Kalman (1960) formalized continuous-time linear \textbf{state-space models} in aerospace engineering -- systems tracking physical positions and velocities via first-order vector differential equations -- establishing the matrix exponential $e^{\mat{A}t}$ as the standard engine for spacecraft trajectory planning and robotic control.

The leap from scalar to matrix exponentials uncovered several classic traps. The most notorious algebraic pitfall is writing $e^{\mat{A}+\mat{B}} = e^{\mat{A}}e^{\mat{B}}$, which fails whenever $\mat{A}$ and $\mat{B}$ do not commute ($\mat{A}\mat{B} \neq \mat{B}\mat{A}$); swapping the order of non-commuting physical operations leads to genuinely different final states. A second trap is numerical: in their classic 1978 paper ``Nineteen Dubious Ways to Compute the Exponential of a Matrix,'' Cleve Moler and Charles Van Loan proved that directly summing the Taylor series $\sum \frac{(\mat{A}t)^k}{k!}$ in floating-point arithmetic is disastrously unstable. Large powers of $\mat{A}$ with alternating signs produce giant intermediate quantities whose subtraction obliterates accuracy -- a numerical breakdown called \textbf{catastrophic cancellation}. When a matrix is defective (possessing repeated eigenvalues that lack a full set of independent eigenvectors), the flow develops polynomial prefactors like $t\,e^{\lambda t}$. In mechanical engineering, these repeated roots correspond to physical \textbf{resonance}, where an external push at a system's natural frequency pumps unbounded energy into an undamped mode.

The matrix exponential provides a continuous-time \textbf{propagator} $\vec{u}(t) = e^{\mat{A}t}\vec{u}_0$ that evolves any initial state forward in time, together with a direct stability test: a system decays to rest if and only if every eigenvalue of $\mat{A}$ has a strictly negative real part. Although real physical dynamics are nonlinear, this linear test remains valid in a tiny neighborhood of an equilibrium point where $\mat{A}$ serves as the \textbf{Jacobian} -- the matrix of all first partial derivatives that provides the best linear approximation to a nonlinear vector field. In 1892, the Russian mathematician Aleksandr Lyapunov showed how to prove stability without solving the differential equation at all, by finding an energy-like scalar function (a \textbf{Lyapunov function}) that steadily decreases along every trajectory.

Two major challenges remain active frontiers today. First, for massive sparse systems with millions of variables (such as power transmission grids, quantum multi-body systems, or chemical reaction networks), explicitly forming or storing the dense $n\times n$ matrix $e^{\mat{A}t}$ is impossible; computing the action $e^{\mat{A}t}\vec{v}$ directly on a state vector using Krylov subspace projections without ever forming the full matrix exponential is an active area of numerical research. Second, \textbf{non-normal matrices} (matrices that do not commute with their transpose, $\mat{A}\mat{A}\T \neq \mat{A}\T\mat{A}$, meaning their eigenvectors are not perpendicular) can trigger colossal \textbf{transient growth}: trajectories can surge thousands of times larger before ultimately decaying, setting off turbulence in fluid flow even when every eigenvalue predicts stability. Analyzing such transient spikes requires the modern theory of \textbf{pseudospectra} and remains an open frontier at the boundary of linear algebra and fluid dynamics. In Chapter~\ref{ch26:what-is-symmetric}, we see how symmetric matrices rule out spirals and guarantee pure exponential decay along perpendicular axes, and in Chapter~\ref{ch27:def-pd}, positive definite matrices provide the energy functions that certify stability.
\end{historybox}

\begin{recoverybox}
\textbf{Cost and intuition recap.} What did solving
$\frac{d\vec{u}}{dt}=\mat{A}\vec{u}$ actually cost, and what is the one idea to
keep?
\begin{itemize}
  \item \textbf{Intuition:} eigenvectors are the directions the flow does not
        mix; along each, the system is just the scalar law $w'=\lambda w$ with
        solution $e^{\lambda t}$. Solve, then superpose.
  \item \textbf{Cost (diagonalizable case):} one eigendecomposition of $\mat{A}$,
        which is $O(n^3)$ work, then \emph{each} value $\vec{u}(t)$ is one
        change of basis, a diagonal scaling, and a change back -- $O(n^2)$ per
        time, for as many times $t$ as you like. You pay the $O(n^3)$ once and
        reuse it forever.
  \item \textbf{Cost (general $\mat{A}$):} the matrix exponential $e^{\mat{A}t}$
        still exists and is computable (scaling-and-squaring of the power series,
        or via Jordan/Schur form); it never needs eigenvectors to be defined.
  \item \textbf{The headline:} the signs of $\operatorname{Re}\lambda_i$ decide
        the fate of \emph{every} trajectory before you compute a single point.
\end{itemize}
\end{recoverybox}

\section*{Summary}
\addcontentsline{toc}{section}{Summary}

\begin{itemize}
  \item A linear system $\frac{d\vec{u}}{dt}=\mat{A}\vec{u}$ is a velocity field;
        its solution is a path flowing along the arrows.
  \item Each eigenpair gives a pure mode $\vec{u}(t)=e^{\lambda t}\vec{x}$. By
        superposition the general (diagonalizable) solution is
        $\vec{u}(t)=\sum_i c_i e^{\lambda_i t}\vec{x}_i$, with $\vec{c}=
        \mat{X}\inv\vec{u}_0$.
  \item \textbf{Stability} is read from the eigenvalues: stable iff every
        $\operatorname{Re}\lambda_i < 0$. The real part is growth/decay; the
        imaginary part is oscillation frequency.
  \item The \textbf{matrix exponential} $e^{\mat{A}t}=\sum_k \mat{A}^kt^k/k!$ is
        the time-$t$ flow map; $\vec{u}(t)=e^{\mat{A}t}\vec{u}_0$ solves the
        system, proved by the same term-by-term differentiation as the scalar
        case.
  \item When $\mat{A}=\mat{X}\mat{\Lambda}\mat{X}\inv$,
        $e^{\mat{A}t}=\mat{X}\,e^{\mat{\Lambda}t}\,\mat{X}\inv$ with
        $e^{\mat{\Lambda}t}=\diag(e^{\lambda_i t})$ -- diagonalization applied to
        continuous time evolution. Defective matrices add polynomial-in-$t$ factors.
  \item Beware: $e^{\mat{A}+\mat{B}}=e^{\mat{A}}e^{\mat{B}}$ only when
        $\mat{A}\mat{B}=\mat{B}\mat{A}$; and continuous stability ($\operatorname{Re}
        \lambda<0$) is not discrete stability ($|\lambda|<1$).
  \item \textbf{Looking forward:} In this chapter, eigenvectors could lean at
        skewed angles and eigenvalues could wander into the complex plane,
        generating spirals. In Chapter~\ref{ch26:what-is-symmetric}, we see
        that when a matrix is symmetric ($\mat{A} = \mat{A}\T$), no spirals can
        occur: every eigenvalue is real and the eigenvectors can always be
        chosen perpendicular, which is the content of the Spectral Theorem.
\end{itemize}

\section{Exercises}
\label{ch25:exercises}

\begin{exercise}[Warmup: scalar fate]
\label{ch25:ex-scalar}
For each, state whether $u(t)\to0$, stays constant, or blows up, and give $u(t)$:
(a) $u'=3u,\ u(0)=2$; (b) $u'=-\tfrac12 u,\ u(0)=5$; (c) $u'=0\cdot u,\ u(0)=7$.
\end{exercise}

\begin{exercise}[Solve a $2\times2$ system]
\label{ch25:ex-solve}
Solve $\frac{d\vec{u}}{dt}=\mat{A}\vec{u}$ with
$\mat{A}=\begin{bmatrix}1&2\\0&-1\end{bmatrix}$ and $\vec{u}(0)=\colvec{1\\1}$.
(Find eigenvalues/eigenvectors, expand the start, write
$\vec{u}(t)$.) Is the system stable? What is the fate of each component as
$t\to\infty$?
\end{exercise}

\begin{exercise}[Read the portrait]
\label{ch25:ex-portrait}
Classify the origin (node/saddle/spiral/center, stable or not) for each:
(a) $\begin{bmatrix}-3&0\\0&-1\end{bmatrix}$;
(b) $\begin{bmatrix}2&0\\0&-1\end{bmatrix}$;
(c) $\begin{bmatrix}-1&-4\\4&-1\end{bmatrix}$;
(d) $\begin{bmatrix}0&3\\-3&0\end{bmatrix}$.
For (c) and (d), find the eigenvalues and identify the real and imaginary parts.
\end{exercise}

\begin{exercise}[Matrix exponential by series]
\label{ch25:ex-series}
Let $\mat{N}=\begin{bmatrix}0&2\\0&0\end{bmatrix}$. Compute $\mat{N}^2$, then sum
the series to find $e^{\mat{N}t}$ exactly. Verify by checking
$\frac{d}{dt}e^{\mat{N}t}=\mat{N}e^{\mat{N}t}$.
\end{exercise}

\begin{exercise}[Diagonalize, then exponentiate]
\label{ch25:ex-diag}
For $\mat{A}=\begin{bmatrix}0&1\\-2&-3\end{bmatrix}$, find $\mat{X}$ and
$\mat{\Lambda}$, then assemble $e^{\mat{A}t}=\mat{X}e^{\mat{\Lambda}t}\mat{X}\inv$
as an explicit $2\times2$ matrix of exponentials. (Eigenvalues are $-1$ and $-2$.)
\end{exercise}

\begin{exercise}[Rotation generator]
\label{ch25:ex-rotation}
Let $\mat{A}=\begin{bmatrix}0&-\omega\\\omega&0\end{bmatrix}$. Using
$\mat{A}^2=-\omega^2\mat{I}$, sum the even and odd parts of the series to show
$e^{\mat{A}t}=\begin{bmatrix}\cos\omega t&-\sin\omega t\\\sin\omega t&\cos\omega t
\end{bmatrix}$. What is the period of the resulting motion?
\end{exercise}

\begin{exercise}[Higher-order ODE as a system]
\label{ch25:ex-secondorder}
The damped oscillator is $y''+3y'+2y=0$. Set $\vec{u}=\colvec{y\\y'}$ and write
the equivalent first-order system $\vec{u}'=\mat{A}\vec{u}$. Find the eigenvalues
of $\mat{A}$ and confirm they are the roots of $\lambda^2+3\lambda+2=0$. Conclude
the general solution $y(t)=c_1 e^{-t}+c_2 e^{-2t}$, and explain why the system is
stable.
\end{exercise}

\begin{exercise}[The commutator obstruction]
\label{ch25:ex-commutator}
Let $\mat{A}=\begin{bmatrix}0&1\\0&0\end{bmatrix}$ and
$\mat{B}=\begin{bmatrix}0&0\\1&0\end{bmatrix}$. Show $\mat{A}\mat{B}\neq
\mat{B}\mat{A}$. Compute $e^{\mat{A}}$ and $e^{\mat{B}}$ (each series is short),
form $e^{\mat{A}}e^{\mat{B}}$, and compare with $e^{\mat{A}+\mat{B}}$ where
$\mat{A}+\mat{B}=\begin{bmatrix}0&1\\1&0\end{bmatrix}$ (this one has eigenvalues
$\pm1$). Confirm they are different, and identify the source of the discrepancy.
\end{exercise}

\begin{exercise}[Trace and determinant of the flow]
\label{ch25:ex-trace}
Prove that $\det\!\bigl(e^{\mat{A}t}\bigr)=e^{t\,\tr\mat{A}}$ for diagonalizable
$\mat{A}$. (Hint: use $e^{\mat{A}t}=\mat{X}e^{\mat{\Lambda}t}\mat{X}\inv$ and that
$\det$ is multiplicative.) Interpret: when is the flow volume-preserving?
\end{exercise}

\begin{exercise}[Proof: stability]
\label{ch25:ex-stability}
Suppose $\mat{A}$ is diagonalizable with all eigenvalues satisfying
$\operatorname{Re}\lambda_i<0$. Using $\vec{u}(t)=\sum_i c_i e^{\lambda_i t}
\vec{x}_i$, prove that $\norm{\vec{u}(t)}\to0$ as $t\to\infty$ for \emph{every}
initial condition. Where exactly does the negativity of the real parts enter?
\end{exercise}

\begin{exercise}[Challenge: the slow mode]
\label{ch25:ex-slowmode}
A stable system has eigenvalues $\lambda_1=-0.2$ and $\lambda_2=-5$ with
eigenvectors $\vec{x}_1,\vec{x}_2$. Argue that for large $t$, \emph{every}
nonzero trajectory $\vec{u}(t)$ becomes nearly parallel to $\vec{x}_1$, regardless
of the start (assuming $c_1\neq0$). Quantify ``large'': roughly how big must $t$
be before the $\vec{x}_2$ contribution is under $1\%$ of the $\vec{x}_1$
contribution, if they start with equal coefficients? (Use $e^{-5t}/e^{-0.2t}=
e^{-4.8t}<0.01$.)
\end{exercise}
